\section{Prospective horizontal side-jet design study}
Relative placement is a design variable, not measured geometry. Nine proposed
arrangements combine three heights within the measured PZT domain with three
nozzle-to-plume distances within the measured positive-centreline ZNMF support.
Selected measured gas-centreline positions are aligned to the existing PZT
fixed-line anchor; nozzle-axis heights differ by the measured gas-centreline offset.
Actual native indices, coordinates, nozzle placements and initial signed velocity
components are exported in the comparison tables. Three adjacent supported native
PZT columns per placement supply 27 representative releases at nominal $d=10\,\mu$m.
These are conditional initialisations in the interaction region, not simulated
unmeasured travel from the atomiser surface.

The separately measured two-component ZNMF mean gas field is prescribed after
rotation/translation of positions and vectors. PZT mean components initialise
particles once and are never added to the gas or repeatedly imposed thereafter.
Finite-Re spherical drag acts on evolving vector slip, with gravity and buoyancy;
constant diameter and no evaporation are assumed. This one-way mean-field model
neglects plume-induced gas changes and jet unsteadiness. It is not experimentally
validated for combined flow. Source cells require four finite vectors; no gas
extrapolation is used and trajectories stop at unsupported boundaries or the
configured 30 ms time horizon. Termination reasons are tabulated. Half-step
comparisons verify position and velocity convergence.

Trajectory Reynolds number uses instantaneous model gas-droplet vector slip.
Supporting $Stk_{D,local}=\tau_{p,Stokes}\lVert\boldsymbol w_g\rVert/D_0$
uses the local nozzle-length gas timescale, not nozzle-exit conditions or the
atomiser frequency. It is not a capture criterion. Deflections and displacements
refer to each release point. No collection target is defined, so no capture
efficiency or optimal arrangement is claimed.

\paragraph{Figure captions.} Trajectory panels: proposed placements over
separately measured mean horizontal gas-velocity contours; curves are integrated
conditional predictions. Velocity panels: horizontal (solid) and vertical (dashed) components for the same releases. Displacement panels use separate horizontal and vertical ordinates. Particle-Re panels use
evolving vector slip; supporting Stokes panels use the explicitly local gas timescale.


% Actual completed study results
\paragraph{Completed results.} All 27 integrations succeeded across the nine
proposed placements. All terminate at the downstream measured-gas boundary,
not at a physical stopping condition. Endpoint deflections compare different
supported elapsed times and are not a ranking of collection performance.
The drag multiplier is unity for $Re_p\leq0.01$ and
$1+0.15Re_p^{0.687}$ otherwise, with $Re_p<800$ enforced.
